\documentclass[10pt]{beamer}
%\usetheme[
%%% option passed to the outer theme
%    progressstyle=fixedCircCnt,   % fixedCircCnt, movingCircCnt (moving is deault)
%] {Feather}
  
% If you want to change the colors of the various elements in the theme, edit and uncomment the following lines



\usetheme{Madrid}
%\usecolortheme{whale}
% Change the bar colors:
%\setbeamercolor{Feather}{fg=red!20,bg=red}
%\setbeamercolor{Feather}{fg=white,bg=white}
%\setbeamercolor{Feather}{fg=purple!25,bg=purple!50}
%\setbeamercolor{Feather}{fg=magenta!25,bg=magenta!65}


% Change the color of the structural elements:
%\setbeamercolor{structure}{fg=red}
%\setbeamercolor{structure}{fg=magenta}
%\setbeamercolor{structure}{fg=purple}

% Change the frame title text color:
%\setbeamercolor{frametitle}{fg=blue}
\setbeamercolor{frametitle}{fg=white}

% Change the normal text color background:
%\setbeamercolor{normal text}{fg=black,bg=gray!10}
%\setbeamercolor{normal text}{fg=black,bg=yellow!10}
%-------------------------------------------------------
% INCLUDE PACKAGES
%-------------------------------------------------------

\usepackage{tkz-graph}
\usepackage[utf8]{inputenc}
\usepackage[brazil]{babel}
\usepackage[T1]{fontenc}
\usepackage{helvet}
%\usepackage{LobsterTwo}
\usepackage{subfigure}



\usepackage{multicol}
%\setlength\columnseprule{0.5pt}
\setlength{\columnsep}{5mm}

\usepackage{amscd}
\usepackage{amsmath}
\usepackage{amsfonts}
\usepackage{amssymb}
\usepackage{amsthm}
\usepackage{latexsym}
\usepackage{eucal}
\usepackage{enumerate}
\usepackage{indentfirst}
\usepackage{booktabs}
\usepackage{bm}

%\usepackage{hyperref}



\newcommand{\La}{\mathscr{L}}
\newcommand{\N}{\mathbb{N}}
\newcommand{\Z}{\mathbb{Z}}
\newcommand{\Q}{\mathbb{Q}}
\newcommand{\R}{\mathbb{R}}
\newcommand{\C}{\mathbb{C}}
\newcommand{\orb}{\mathcal{O}}

%Disjoint Union
\newcommand{\cupdot}{\mathbin{\mathaccent\cdot\cup}}


%%%%%%%%%%%%%%%%%%%%%%%%%%%
%OPERADORES MATEMÁTICOS
% \sen , \arcsen , etc.
%%%%%%%%%%%%%%%%%%%%%%%%%%%
\DeclareMathOperator{\sen}{sen}
\DeclareMathOperator{\arcsen}{arcsen}
\DeclareMathOperator{\senh}{senh}
\DeclareMathOperator{\arccot}{arccot}
\DeclareMathOperator{\arcsec}{arcsec}
\DeclareMathOperator{\arccsc}{arccsc}
\DeclareMathOperator{\traco}{traço}
\DeclareMathOperator{\mmc}{mmc}
\DeclareMathOperator{\mdc}{mdc}
\DeclareMathOperator{\image}{Im}
\DeclareMathOperator{\vspan}{span}
\DeclareMathOperator{\e}{\mathrm{e}}
\DeclareMathOperator{\tr}{tr}
\DeclareMathOperator{\rot}{rot}
\DeclareMathOperator{\grad}{grad}
\DeclareMathOperator{\divv}{div}


\newcommand{\dlim}{\displaystyle\lim}
\newcommand{\dint}{\displaystyle\int}
\newcommand{\doint}{\displaystyle\oint}
\newcommand{\dsum}{\displaystyle\sum}
\newcommand{\lsum}{\sum\limits}


\newtheorem{def1}[section]{Definição}
\newtheorem{teo}[section]{Teorema}

%----------------------------------------------------------------------------------------
%	TITLE PAGE
%----------------------------------------------------------------------------------------



\begin{document}


\title[Mecânica Clássica e Quântica]{Mecânica Clássica e Quântica} % The short title appears at the bottom of every slide, the full title is only on the title page


\author{  Bianca Bonfim\\
Pedro Cruz \\
Richard Francis\\

} % Your name
\institute[MCCT] % Your institution as it will appear on the bottom of every slide, may be shorthand to save space
{
}

%\date{\today} % Date, can be changed to a custom date
\begin{frame}
	\maketitle

\end{frame}


\begin{frame}{Sumário}
\tableofcontents
\end{frame}

%---------------------------------------
\section{Problema}
\begin{frame}{Problema de Dois Corpos}
    Considere o problema de interação gravitacional entre dois corpos esféricos que podem
ser considerados como partículas pontuais com massa constante $m_1$ e $m_2$. 

Utilizando a Mecânica Newtoniana e conhecendo as condições iniciais $\vec{r_1}(t = 0) = \vec{r_1}(0)$,
$\vec{v_1}(t = 0) =\vec{v_1}(0)$ e $\vec{r_2}(t = 0) =\vec{r_2}(0)$, $\vec{v_2}(t = 0) = \vec{v_2}(0)$, devemos encontrar a posição de ambos corpos em cada instante de tempo $t$.
\end{frame}
\begin{frame}{Problema de dois corpos}
   Considere duas particulas pontuais de massas constantes $m_1$ e $m_2$, nas posicoes
$\bm r_1(t)$ e $\bm r_2(t)$. A interacao gravitacional entre elas é dada por:
\begin{equation}
\bm F_{12}
=G\frac{m_1m_2}{|\bm r_2-\bm r_1|^3}(\bm r_2-\bm r_1),
\qquad
\bm F_{21}=-\bm F_{12}.
\end{equation}

Definindo
\begin{equation}
\bm r=\bm r_2-\bm r_1,\qquad r=|\bm r|,
\end{equation}
o potencial gravitacional é
\begin{equation}
U(\bm r)=-\frac{Gm_1m_2}{r}.
\end{equation}

Tambem definimos a massa total, a massa reduzida e o centro de massa:
\begin{equation}
M=m_1+m_2,\qquad
\mu=\frac{m_1m_2}{m_1+m_2},\qquad
\bm R=\frac{m_1\bm r_1+m_2\bm r_2}{M}.
\end{equation}

\end{frame}
\section{Mecânica Newtoniana}
\begin{frame}{Mecânica Newtoniana}
    As equacoes de Newton sao
\begin{equation}
m_1\ddot{\bm r}_1
=G\frac{m_1m_2}{r^3}\bm r,
\qquad
m_2\ddot{\bm r}_2
=-G\frac{m_1m_2}{r^3}\bm r.
\end{equation}

Somando as duas equações:
\begin{equation}
M\ddot{\bm R}=0.
\end{equation}

Portanto, o centro de massa realiza movimento retilíneo uniforme:
\begin{equation}
\bm R(t)=\bm R_0+\bm V_{CM}t,
\end{equation}
com
\begin{equation}
\bm R_0=\frac{m_1\bm r_1(0)+m_2\bm r_2(0)}{M},
\qquad
\bm V_{CM}=\frac{m_1\bm v_1(0)+m_2\bm v_2(0)}{M}.
\end{equation}
\end{frame}

\begin{frame}{Mecânica Newtoniana}
    Subtraindo as acelerações, obtemos a equação relativa:
\begin{equation}
\ddot{\bm r}=-\frac{GM}{r^3}\bm r.
\end{equation}

Assim, o problema de dois corpos se reduz ao movimento de uma particula fictícia
de massa reduzida $\mu$ sob o potencial central:
\begin{equation}
U(r)=-\frac{GM\mu}{r}.
\end{equation}

As condicoes iniciais relativas são:
\begin{equation*}
\bm r_0=\bm r_2(0)-\bm r_1(0),
\qquad
\bm v_0=\bm v_2(0)-\bm v_1(0).
\end{equation*}
\end{frame}

\begin{frame}{Mecânica Newtoniana}
    Conhecida a solucao relativa $\bm r(t)$, as posicoes reais dos corpos sao
\begin{equation}
\boxed{
\bm r_1(t)=\bm R_0+\bm V_{CM}t-\frac{m_2}{M}\bm r(t)
}
\end{equation}
e
\begin{equation}
\boxed{
\bm r_2(t)=\bm R_0+\bm V_{CM}t+\frac{m_1}{M}\bm r(t)
}.
\end{equation}
\end{frame}

\section{Mecânica Lagrangeana}
\begin{frame}{Mecânica Lagrangeana}
Usaremos como coordenadas generalizadas as coordenadas cartesianas
$r_1 = (x_1, y_1, z_1)$ e $r_2 = (x_2, y_2, z_2)$. Como $L=T-U$ então a Lagrangiana do sistema é:
\begin{equation*}
L=
\frac{1}{2}m_1\dot{\bm r}_1^{\,2}
+\frac{1}{2}m_2\dot{\bm r}_2^{\,2}
+\frac{Gm_1m_2}{r}.
\end{equation*}

    Usando as coordenadas de centro de massa $\bm R$ e $\bm r$, obtém-se
    \begin{equation*}
T=
\frac{1}{2}M\dot{\bm R}^{\,2}
+\frac{1}{2}\mu\dot{\bm r}^{\,2}.
\end{equation*}

Onde $\mu=\frac{m_1 m_2}{m_1+m_2}$ e $M=m_1+m_2$.
Assim,
\begin{equation*}
L(\dot{\bm R},\bm r,\dot{\bm r})
=
\frac{1}{2}M\dot{\bm R}^{\,2}
+\frac{1}{2}\mu\dot{\bm r}^{\,2}
+\frac{GM\mu}{r}.
\end{equation*}
\end{frame}

\begin{frame}{Mecânica Lagrangeana}
Pelas equações de Euler-Lagrange, obtém-se para $\bm R$
\begin{equation*}
\frac{d}{dt}
\left(
\frac{\partial L}{\partial \dot {\bm R}}
\right) =
\frac{\partial L}{\partial \bm R}=0.
\end{equation*}
E portanto,
$$M \ddot{\bm R}=0$$
A mesma equação obtida por Newton.

Para a coordenada relativa:
\begin{equation*}
\frac{\partial L}{\partial \bm r}
=GM\mu\frac{\partial}{\partial \bm r}
\left(\frac{1}{r}\right)
=-GM\mu\frac{\bm r}{r^3}.
\end{equation*}

Assim,
\begin{equation*}
\mu\ddot{\bm r}+GM\mu\frac{\bm r}{r^3}=0,
\end{equation*}
ou
\begin{equation*}
\boxed{
\ddot{\bm r}=-\frac{GM}{r^3}\bm r
}.
\end{equation*}
\end{frame}

\section{Mecânica Hamiltoniana}
\begin{frame}{Mecânica Hamiltoniana}
    Os momentos conjugados são:
\begin{equation}
\bm P=\frac{\partial L}{\partial \dot{\bm R}}=M\dot{\bm R},
\qquad
\bm p=\frac{\partial L}{\partial \dot{\bm r}}=\mu\dot{\bm r}.
\end{equation}

A Hamiltoniana e obtida por transformacao de Legendre:
\begin{equation}
H=\bm P\cdot\dot{\bm R}+\bm p\cdot\dot{\bm r}-L.
\end{equation}

Como
\begin{equation}
\dot{\bm R}=\frac{\bm P}{M},
\qquad
\dot{\bm r}=\frac{\bm p}{\mu},
\end{equation}
segue que
\begin{equation}
\boxed{
H(\bm R,\bm P,\bm r,\bm p)
=\frac{P^2}{2M}+\frac{p^2}{2\mu}-\frac{GM\mu}{r}
}.
\end{equation}

As equações de Hamilton são:
\begin{equation}
\dot q_i=\frac{\partial H}{\partial p_i},
\qquad
\dot p_i=-\frac{\partial H}{\partial q_i}.
\end{equation}

Para o centro de massa:
\begin{equation}
\dot{\bm R}=\frac{\bm P}{M},
\qquad
\dot{\bm P}=0.
\end{equation}
\end{frame}

\begin{frame}{Mecânica Hamiltoniana}
    
Logo,
\begin{equation}
\bm P=\text{constante},
\qquad
\bm R(t)=\bm R_0+\frac{\bm P_0}{M}t.
\end{equation}

Para o movimento relativo:
\begin{equation}
\dot{\bm r}=\frac{\bm p}{\mu},
\end{equation}
\begin{equation}
\dot{\bm p}
=-\frac{\partial}{\partial \bm r}
\left(-\frac{GM\mu}{r}\right)
=-GM\mu\frac{\bm r}{r^3}.
\end{equation}

Como $\bm p=\mu\dot{\bm r}$, resulta
\begin{equation}
\boxed{
\ddot{\bm r}=-\frac{GM}{r^3}\bm r
}.
\end{equation}

Assim, o formalismo Hamiltoniano tambem reproduz a solução de Kepler.

\end{frame}

\begin{frame}{Caso de massas variáveis}
    Agora considere
$m_1(t)=m_{10}e^{-t/\tau_1}$ e
$m_2(t)=m_{20}e^{-t/\tau_2},$
com $m_{10}$, $m_{20}$, $\tau_1$ e $\tau_2$ constantes positivas.

De $F=d/dt(m\bm v)$, e $\bm r = |r_2 - r_1|$ temos:

\begin{equation}
\frac{d}{dt}(m_1\dot{\bm r}_1)
=G\frac{m_1m_2}{r^3}\bm r,
\end{equation}
\begin{equation}
\frac{d}{dt}(m_2\dot{\bm r}_2)
=-G\frac{m_1m_2}{r^3}\bm r.
\end{equation}

Somando as equações, chegamos em:
\begin{equation}
\frac{d}{dt}
(m_1\bm v_1+m_2\bm v_2)=0.
\end{equation}

Portanto, nessa convenção, o momento linear total ainda se conserva. Entretanto, a
energia mecânica usual nao se conserva, pois o potencial depende explicitamente do tempo.
\end{frame}
\section{Simulações Computacionais}
\begin{frame}{Simulações Computacionais}
    A simulação foi realizada em Python considerando os seguintes casos:
    
\begin{enumerate}
    \item $m_1=m_2$, particulas proximas;
    \item $m_1=m_2$, particulas nao tao proximas;
    \item $m_1=333000m_2$, particulas proximas;
    \item $m_1=333000m_2$, particulas nao tao proximas.
\end{enumerate}

Em cada caso, são consideradas duas condições iniciais:
\begin{itemize}
    \item partículas inicialmente em repouso;
    \item partículas com velocidade inicial diferente de zero.
\end{itemize}
\end{frame}

\begin{frame}{Conclusão}
    Newton, Lagrange e Hamilton levam à mesma solução, mas organizam o problema de
formas diferentes: forças, ação e espaço de fase, respectivamente. As conservações de
momento linear, energia e momento angular seguem da ausência de forças externas, da
independência temporal da Lagrangiana e da simetria rotacional do potencial central.
\end{frame}

\begin{frame}{Referências}
 \begin{thebibliography}{99}

\bibitem{}
BENITEZ, Gustavo.
\textit{Notas de Aula de Mecânica Clássica}.
Universidade Federal Fluminense (UFF), 2026.

\end{thebibliography}
\end{frame}

\end{document}